\documentclass[12pt]{article}
\usepackage{fontspec}\usepackage{xeCJK}\usepackage{geometry}
\usepackage{fancyhdr}\usepackage{titlesec}\usepackage{caption}
\usepackage{graphicx}\usepackage{amsmath}\usepackage{booktabs}
\usepackage{tikz}\usetikzlibrary{positioning,arrows.meta,fit,calc,backgrounds}
\usepackage[RPvoltages]{circuitikz}
\usepackage{listings}\usepackage{xcolor}

% Rust 代码清单样式（仅 ASCII，避免 CJK 进 listing）
\definecolor{cmt}{RGB}{110,120,130}
\definecolor{kw}{RGB}{20,70,160}
\definecolor{str}{RGB}{160,60,40}
\lstdefinestyle{rust}{
  basicstyle=\ttfamily\footnotesize, keywordstyle=\color{kw}\bfseries,
  commentstyle=\color{cmt}\itshape, stringstyle=\color{str},
  numbers=left, numberstyle=\tiny\color{cmt}, numbersep=6pt,
  showstringspaces=false, breaklines=true, columns=fullflexible,
  frame=single, framesep=4pt, xleftmargin=14pt, aboveskip=6pt, belowskip=2pt,
  morekeywords={let,mut,fn,for,in,if,else,while,loop,struct,impl,pub,const,
    as,return,self,Self,f32,i32,u32,i16,u16,usize,i64,u64}}

\geometry{a4paper, portrait, top=3cm, bottom=2cm, left=2.5cm, right=2.5cm,
  headheight=0pt, headsep=0pt, footskip=1.2cm}

\setmainfont{Times New Roman}
\setCJKmainfont{STSong} % macOS 无 SimSun；STSong 同为宋体。Windows/Overleaf 可改回 SimSun
\xeCJKsetup{AutoFakeBold=2.5}

\newcommand{\sanhao}{\fontsize{16pt}{22pt}\selectfont}
\newcommand{\sihao}{\fontsize{14pt}{22pt}\selectfont}
\newcommand{\xiaosi}{\fontsize{12pt}{22pt}\selectfont}
\newcommand{\wuhao}{\fontsize{10.5pt}{15pt}\selectfont}

\titleformat{\section}{\sanhao\bfseries}{\thesection}{1em}{}
\titleformat{\subsection}{\sihao\bfseries}{\thesubsection}{1em}{}
\titleformat{\subsubsection}{\xiaosi\bfseries}{\thesubsubsection}{1em}{}
\titlespacing{\section}{0pt}{12pt}{6pt}
\titlespacing{\subsection}{0pt}{6pt}{6pt}

\DeclareCaptionFont{wuhao}{\wuhao}
\captionsetup{font=wuhao, labelsep=quad}
\renewcommand{\figurename}{图}
\renewcommand{\tablename}{表}

% 浮动排布：放宽比例限制，避免表格聚集成难看的纯浮动页
\renewcommand{\textfraction}{0.05}
\renewcommand{\topfraction}{0.95}
\renewcommand{\bottomfraction}{0.95}
\renewcommand{\floatpagefraction}{0.85}

\pagestyle{fancy}\fancyhf{}
\renewcommand{\headrulewidth}{0pt}
\fancyfoot[R]{\wuhao\thepage}


\begin{document}
\xiaosi

\newpage
\setcounter{page}{1}

% ===== 摘要 =====
\begin{center}
  {\sanhao\bfseries 摘\quad 要}
\end{center}
\noindent 本作品完成了一个无源交流电流表及无线读表器，该表由变压器，电流互感器，电磁耦合供电模块，TI MSPM0G507 单片机四部分构成。

\vspace{6pt}
\noindent{\xiaosi\bfseries 关键词：}无线读表器；电流互感器；

\newpage
\setcounter{page}{1}

% ===== 正文 =====

% TODO 硬件：系统方案、取电前端、传感通道、放电回路、读表器硬件

\section{程序设计}
\subsection{软件总体结构}
软件以测量周期为单位循环执行，流程见图~\ref{fig:flow}。每周期依次完成探测帧采集与量程判定、
测量帧采集、有效值计算、标定换算、显示与无线发送。采样由定时器触发 ADC、DMA 成帧，采样间隔
由时钟决定。主要参数见表~\ref{tab:param}。

\begin{figure}[htbp]\centering
\begin{tikzpicture}[
  font=\wuhao,
  blk/.style={draw, rounded corners=2pt, align=center, minimum height=10mm,
    inner sep=2pt, fill=black!3, text width=17mm},
  >={Stealth[length=2mm]}, line width=0.4pt, node distance=3mm]
  \node[blk] (probe) {探测帧\\$\times1$ 直采\\160 点/40\,ms};
  \node[blk, right=of probe] (pick) {量程判定\\$G$、DAC 码};
  \node[blk, right=of pick] (main) {测量帧\\4\,ksps\\800 点};
  \node[blk, right=of main] (rms) {加权有效值\\（LSB）};
  \node[blk, right=of rms] (cal) {分档标定\\（A）};
  \node[blk, right=of cal] (out) {OLED 显示\\无线发送};
  \draw[->] (probe) -- (pick);  \draw[->] (pick) -- (main);
  \draw[->] (main) -- (rms);    \draw[->] (rms) -- (cal);
  \draw[->] (cal) -- (out);
  \draw[->] (out.south) -- ++(0,-3mm) -| (probe.south);
\end{tikzpicture}
  \caption{测量周期流程}
  \label{fig:flow}
\end{figure}

\begin{table}[htbp]\centering
\caption{采样与测量参数}
\label{tab:param}
\wuhao
\begin{tabular}{ll}
\toprule
参数 & 取值\\
\midrule
采样频率 $f_s$      & 4\,kHz\\
帧长 $N$            & 800 点（200\,ms，10 个工频周期）\\
探测帧长            & 160 点（40\,ms）\\
ADC                 & 12 位，参考片内 2.5\,V 基准，满量程 $0\text{--}4095$\\
DAC                 & 12 位，与 ADC 共用同一基准\\
增益档位            & $\times1$、$\times2$、$\times4$、$\times8$、$\times16$、$\times32$\\
窗函数              & 周期型 Hann\\
\bottomrule
\end{tabular}
\end{table}

\subsection{交流电流有效值的数字测量}
交流电流有效值定义为 $I=\sqrt{\frac1T\int_0^T i^2(t)\,\mathrm{d}t}$。对等间隔采样序列 $x[n]$，
取整数枢轴 $p$、窗函数 $w[n]=\tfrac12\bigl(1-\cos\frac{2\pi n}{N}\bigr)$，加权直流与加权有效值为
\begin{equation}
  \delta=\frac{\sum_n w[n]\,(x[n]-p)}{\sum_n w[n]},\qquad
  X_{\text{rms}}=\sqrt{\frac{\sum_n w[n]\,\bigl(x[n]-p-\delta\bigr)^2}{\sum_n w[n]}}\,.
\end{equation}
帧长取工频周期的整数倍，$N=f_s M/f=800$（$M=10$，$f=50$\,Hz）。误差按下列各项估算，汇总于
表~\ref{tab:budget}。

\textbf{非同步采样误差}\quad 工频与采样时钟均存在偏差，实际帧长不严格为整周期，平方和中残留
二倍频分量，其幅度等于窗序列频响在 $2f$ 处的取值。本设计 $2f$ 对应第 20 号频点：矩形窗该处
旁瓣约 $-35$\,dB，按 49.66\,Hz 计对应有效值误差 0.33\,\%；Hann 窗低于 $-95$\,dB，该项可忽略。
$h$ 次谐波对应第 $20h$ 号频点，抑制更深。因此本设计不需锁频与过零检测。

\textbf{量化误差}\quad 量化噪声有效值 $q/\sqrt{12}=0.29$\,LSB。0.1\,A 在 $\times32$ 档的读数约
233\,LSB，量化项为 0.12\,\%；不加放大时约 7\,LSB，量化项达 4\,\%。

\textbf{噪声偏置}\quad 噪声与信号正交叠加，读数为 $\sqrt{X^2+\sigma^2}$，相对偏差
$\sigma^2/2X^2$，为单边误差，多帧平均不能消除。该项由前端纹波决定，需实测 $\sigma$ 后确定。

\begin{table}[htbp]\centering
\caption{有效值测量误差预算（0.1\,A，$\times32$ 档，指标 0.5\,\%）}
\label{tab:budget}
\wuhao
\begin{tabular}{lll}
\toprule
误差项 & 量值 & 来源\\
\midrule
非同步采样（Hann 窗） & $<0.01\,\%$ & 分析值，矩形窗为 $0.33\,\%$\\
量化                  & $0.12\,\%$  & $0.29$\,LSB / 233\,LSB\\
单帧噪声底            & $0.043\,\%$ & 分析值，加窗前 $0.035\,\%$\\
噪声偏置 $\sigma^2/2X^2$ & $\sigma=1$\,LSB 时 $0.06\,\%$；$\sigma=3$\,LSB 时 $0.53\,\%$ & 待实测 $\sigma$\\
标定表插值残差        & 待实测 & 由标定点间距决定\\
\bottomrule
\end{tabular}
\end{table}

\subsection{自动量程与偏置}
$0.1\text{--}2$\,A 量程比为 20:1，指标 0.5\,\% 在下端为 0.5\,mA，故采用逐周期自动量程。片内运放
配置为非反相可编程增益放大器，电阻梯底部接片内 12 位 DAC，输出电压、偏置码与选档条件为
\begin{equation}
  V_{\text{out}} = G\,V_{\text{in}} + (1-G)\,V_{\text{dac}},\qquad
  D=\frac{G\,m-C}{G-1},\qquad
  \frac{G\,p}{2}\le 0.65\min(d,\,4095-d),
\end{equation}
式中 $m$、$p$ 为探测帧测得的输入直流与峰峰值，$C=2047$ 为窗口中心，$d$ 为该档输出直流，取满足
条件的最大档。DAC 使放大以输入直流为支点；直流灵敏度 $\mathrm{d}V_{\text{out}}/\mathrm{d}V_{\text{in,dc}}=G$，
$\times32$ 档上 50\,mV 失配即使输出偏离中心 1.6\,V。$\times1$ 档由 ADC 直接读取输入，用于输入
直流过低、各 PGA 档均无法容纳的情形。换档设 25\,\%/75\,\% 回差。

\subsection{标定}
标定表为"加权有效值（LSB）$\to$ 电流（A）"的分段线性表，每档一张，读数按
$I=y_0+(y_1-y_0)(X-x_0)/(x_1-x_0)$ 插值，表外按端段外推。分档建表的原因是片内电阻梯各档实际
比值对标称 $2/4/8/16/32$ 的偏离与 0.5\,\% 指标同量级，且互感器在低磁通下比差随电流变化，单一
比例系数无法覆盖。标定数据由台架比对获得，填表方法见测试章节。

\subsection{通信协议}
无线链路为透传串口，不提供分包与长度字段，故采用面向行的 ASCII 协议：以换行为自同步定界符，
接收端在任意位置接入或丢失字节后都能在下一个换行重新对齐；不合格式的行整行丢弃。链路速率
115200\,bps，一次读数的两行帧约 110 字节，线上时间不足 10\,ms。

题目 2(2) 规定两种读取模式下均不得操作电流表，故读数由读表器请求产生，电流表不主动上报测量
值。帧格式见表~\ref{tab:frame}。

\begin{table}[htbp]\centering
\caption{帧格式}
\label{tab:frame}
\wuhao
\begin{tabular}{llp{6.2cm}}
\toprule
方向 & 帧 & 说明\\
\midrule
读表器 $\to$ 电流表 & \texttt{MEAS} & 广播，任一在线电流表应答\\
读表器 $\to$ 电流表 & \texttt{MEAS,ADDR=$n$} & 仅编码开关为 $n$ 的电流表应答\\
电流表 $\to$ 读表器 & \texttt{METER\_TEST,ADDR=$n$,CURRENT\_MA=$m$,STATUS=$s$}
  & 读数帧。$s\in\{$\texttt{NORMAL},\texttt{LOW},\texttt{HIGH}$\}$，按 0.2\,A/2\,A 判定\\
电流表 $\to$ 读表器 & \texttt{METER\_CAL,ADDR=$n$,RMS=$x$,GAIN=$g$,FLAG=$f$,\dots}
  & 同次读数的原始加权有效值、档位与质量标志，供标定与诊断\\
电流表 $\to$ 读表器 & \texttt{IMHERE,ADDR=$n$} & 保活心跳，空闲时每 2\,s 一次，不含读数\\
\bottomrule
\end{tabular}
\end{table}

命令带地址而非仅用广播，是因为一台电流表在不同时刻由编码开关赋予不同地址，仅凭连接无法确定
对端身份；带地址后请求与应答一一对应，不会把某地址的应答记到另一地址名下。

心跳不参与测量，不唤醒任何模拟电路，其作用有二：编码开关改变时旧地址的心跳停止、新地址出现，
读表器据此更新在场地址；负载断开后电流表随之失电，心跳消失，即为题目 2(3) 的离线判据，较等待
一次请求超时更快。

读数不可信时（基准未就绪、输入越出转换窗口、采集帧未完成）电流表只发 \texttt{METER\_CAL} 并在
\texttt{FLAG} 中给出原因，不发读数帧，避免把不可信的数值以正常状态送出。

\subsection{读表器软件}
读表器周期扫描（8\,s 扫描、12\,s 间歇）并对搜索到的电流表建立连接，连接建立后不再持续收数：
读数由请求产生，链路空闲是常态，故链路健康仅由连接状态判定，不以数据超时判定。

基本模式为一对一手动读取，操作读表器选定地址后发出 \texttt{MEAS,ADDR=$n$}，收到读数帧即显示
并记录。自动模式一键启动，每 2 分钟对当前在场（心跳可闻）的地址依次发出请求；同一连接上同时
只允许一个未完成请求，故轮次内串行执行。

在场与离线由心跳判定：收到该地址任一帧即刷新其在场时刻，连续 7\,s（三个心跳周期）无帧则判为
离线并给出指示。离线只作指示，不写入记录——记录用于保存读数，而缺席不是读数。请求超时同样
不写入记录，只提示本次未读到，且不据此判定离线：正在心跳的电流表只是本次未及应答，并非不在。

报警按读数判定，低于 0.2\,A 为低限、高于 2\,A 为超限，每一次越限读数均给出声音与振动提示，不
以状态跳变为条件——读数稀疏，同一状态的相邻两次读数是两个独立事件。

记录存入手机数据库，含读取时间、电流表地址、电流值与状态，掉电不丢失，容量 1000 条，并按地址
绘制电流趋势曲线。

% TODO 测试：\section{测试方案与测试结果}\label{sec:test}

\appendix
% TODO 附录 A：前端电路原理图

\section{附录 B\quad 有效值计算源码}
一帧 800 点的 Hann 加权真有效值。均值与平方和用同一个窗加权，并分两趟：合成一趟要用
$\overline{wx^{2}}-\bar{x}^{2}$，是两个 $4.2\times10^{6}$ 量级的数相减去取 $10^{2}$ 量级的结果，
单精度浮点的尾数不够这次抵消。

\begin{lstlisting}[style=rust,caption={Hann 加权真有效值}]
// Pass 1: Hann-weighted DC, as an offset from the integer pivot.
let (mut sum_w, mut sum_wd) = (0.0f32, 0.0f32);
for (&x, &w) in buf.iter().zip(HANN.iter()) {
    sum_w += w;
    sum_wd += w * ((x & 0xFFF) as i32 - pivot) as f32;
}
let offset = sum_wd / sum_w;

// Pass 2: weighted mean square of the residual.
let mut sum_wd2 = 0.0f32;
for (&x, &w) in buf.iter().zip(HANN.iter()) {
    let d = ((x & 0xFFF) as i32 - pivot) as f32 - offset;
    sum_wd2 += w * d * d;
}
sqrtf(sum_wd2 / sum_w) // rms, in ADC LSB
\end{lstlisting}

\end{document}
